\section{Introduction}

\noindent Approximation algorithms and fixed parameter-tractabililty (FPT) are two ways to cope with NP-hard problems. Recently, there have been many works that unite the two by obtaining approximation algorithms for NP-Hard problems that run in FPT time. Examples include {\sc Vertex Coloring}\cite{DHK05,Mar08}, {\sc $k$-Path Deletion}\cite{Lee19}, {\sc Vertex Cycle Packing}\cite{Lok+17}, {\sc Flow Time Scheduling}\cite{Wie18}, {\sc Max $k$-Vertex Cover} in $d$-uniform hypergraphs \cite{SF17,Man19}, {\sc $k$-Means} and {\sc $k$-Median} \cite{LiS16,BPR+17,KMN+04,ANSW20,CGTS02,CGK+19}. 
On the other hand, there are also various developments in FPT hardness of approximation, for example, for {\sc $k$-Biclique} \cite{Lin18}, \kclique{} \cite{CCK+17, Lin21, LRSW22, KK22, LRSW23, CFL+23}, \ksetcover{} \cite{CCK+17,CL19,KLM19,Lin19,KN21,LRSW23a} and so on.
We refer to the survey by Feldmann, Karthik, Lee, and Manurangsi~\cite{FKLM20} for a more comprehensive list of both FPT approximation algorithms and FPT hardness of approximation results.

However, it is worth pointing out that the techniques used in proving FPT hardness of approximation results have been rather problem-specific, and there remain other basic problems for which even constant inapproximability is unknown. This situation is due to the lack of a PCP-like theorem in the FPT world. In classical polynomial time hardness of approximation, the PCP theorem is a fundamental result that leads to a plethora of other results via reductions. The analog of the PCP theorem in the FPT regime was explicitly proposed by Lokshtanov et al~\cite{LRSZ20} and named \emph{Parameterized Inapproximability Hypothesis (PIH)}. The PIH states that there is an absolute constant $\varepsilon > 0$ such that for a 2CSP on $k$ variables with alphabet size $n$, there is no algorithm running in time $f(k)\cdot n^{O(1)}$ that can distinguish a satisfiable instance from one where at most $(1-\varepsilon)$ fraction of constraints can be simultaneously satisfied.

Analogous to the PCP theorem, PIH not only implies many FPT time inapproximability results including \kclique{}, \ksetcover{}, \kexact{}, and {\sc $k$-SetPacking}, but also unifies the inapproximability landscape by serving as a versatile problem-agnostic starting point.  Previously, PIH was known to be implied by Gap-ETH \cite{DM18, CFM17}, a strong assumption with an inherent gap in it. Establishing PIH under a gap-free hypothesis has been a significant open problem in parameterized complexity. Very recently, following the posting of this work, PIH was established under ETH \cite{GLR+23} using the technique of parallelization. The authors first reduce 3SAT to a vectorized problem called Vector-Valued CSP, and then design parallel probabilistically checkable proofs to verify its satisfiability. However, the proof in \cite{GLR+23} fails to get any inapproximability of 2CSP under $\wone \neq \fpt$, since it is not known how to reduce a $\wone$-complete problem, say $\kclique{}$, to their Vector-Valued CSP. In fact, it was pointed out in \cite{GLR+23} that Vector-Valued CSP is likely to be $\mone$-complete and thus not likely to be $\wone$-hard, where $\mone$ is an intermediate complexity class between $\fpt$ and $\wone$ \cite{fellows2003blow, chen2007isomorphism}. Proving PIH under $\wone \neq \fpt$, therefore, still remains an important open problem.

\begin{comment}
 So far, however, PIH has only been established under very strong assumptions such as Gap-ETH \cite{DM18,CFM17}. Even ETH, a stronger assumption than $\wone\neq\fpt$ \cite{CHKX06}, is not known to imply PIH. 
\vnote{Edit this sentence - could say that recently, following the posting of this work, PIH was shown under ETH}
Establishing PIH under $\wone \neq \fpt$ remains a significant major open problem in parameterized complexity.
\end{comment}

In this work, we prove a list version of the PIH, which we call \emph{Baby PIH}, under the minimal hypothesis $\wone \neq \fpt$. In Baby PIH, we study $r$-list assignments to the underlying 2CSP instance. An $r$-list assignment has a list $L(u)$ of at most $r$ values to each variable $u$, and a constraint between variables $u,v$ is said to be satisfied by the $r$-list assignments if there is at least one pair of values $x \in L(u), y \in L(v)$ such that $(x,y)$ satisfies the constraint between the variables $u,v$. 

\begin{definition}[Baby PIH]
For all constants $r, C \ge 1$ and every computable function $f$, there is no algorithm running in time $f(k)\cdot n^C$ that can distinguish between the following cases, on input a 2CSP instance on $k$ variables with alphabet size $n$. 
    \begin{itemize}
        \item (Completeness) There is an assignment satisfying all the constraints. 
        \item (Soundness) No $r$-list assignment satisfies all the constraints. 
    \end{itemize}
\end{definition}

The Baby PIH can be compared with the ``Baby PCP theorem,'' which was established via a purely combinatorial gap theorem in a remarkable recent work of Barto and Kozik~\cite{BK22} (who also coined the term Baby PCP). 
The Baby PCP theorem asserts the NP-completeness of distinguishing satisfiable instances of a 2CSP from those which lack a satisfying $r$-list assignment, for all constants $r > 1$.   Baby PCP differs from Baby PIH in the sense that it concerns 2CSP on $n$ variables with constant alphabet size.  The concept of Baby PCP is itself not new and was studied in the early PCP days under the guise of a certain minimization version of Label Cover~\cite{ABSS97} (see also \cite{DS04}).\footnote{Label Cover is stronger than Baby PCP as stated above, as it imposes that the 2CSP relations are functions (usually referred to as projection property in PCP parlance). However, Barto and Kozik's version also has this projection property.} 

The term ``Baby'' stems from the fact that after reducing the soundness parameter to an arbitrarily small positive constant via parallel repetition, one can deduce the Baby PCP theorem from the PCP theorem.  This strategy also works in the PIH setting. Namely, we can use the canonical ``clause-variable'' construction to build a 2CSP instance with projection property, then apply parallel repetition \cite{Rao11} to amplify the soundness to below $\frac{1}{r^2}$. Suppose a 2CSP instance is $r$-list satisfiable, then by randomly picking a value from each variable's list, we can conclude there is a single assignment that satisfies at least $\frac{1}{r^2}$ fraction of constraints. Therefore, being unable to approximate 2CSP within a $\frac{1}{r^2}$ factor implies that it is hard to distinguish whether a instance is satisfiable or not $r$-list satisfiable. In other words, PIH implies Baby PIH. Thus establishing the Baby PIH is a necessary step towards proving the PIH itself, and we also believe that it is a valuable intermediate step. 

In this work, we prove Baby PIH under the minimal complexity assumption $\wone\neq\fpt$. 

\begin{theorem}
\label{thm:main-intro}
    Assuming $\wone\neq\fpt$, Baby PIH is true.
\end{theorem}

Our proof of Theorem~\ref{thm:main-intro} is combinatorial and follows the framework of Barto and Kozik's proof of the Baby PCP theorem. Specifically, given a 2CSP instance with variable set $X$ and a list size $r$, we choose large enough integers $a,b$ depending on $r$, and construct a bipartite direct product 2CSP instance, whose variable set is $\binom{X}{a}\cup\binom{X}{b}$. Given that the product instance is $r$-list satisfiable, we can repeatedly choose smaller integers $a' \le a, b' \le b$ and extract assignments for the instance with variable set $\binom{X}{a'}\cup \binom{X}{b'}$. The new assignments still list satisfy the smaller instance, but the size of each list on one side is decreased by 1, which helps us to do induction.

We highlight that although this proof strategy looks simple, there are basic obstacles to employ it in the Baby PIH setting (compared to the Baby PCP setting). In the PCP world, the alphabet size $|\Sigma|$ is at most some constant. Thus, to extract assignments for the smaller instance, it is affordable to pick $a,b$ large enough depending on $|\Sigma|$. The running time of there reduction is therefore $|X|^{O_{r,|\Sigma|}(1)}$. However in the PIH case, $|\Sigma|$ can be as large as the input size, and the running time of the reduction can only be $f(|X|)\cdot |\Sigma|^{O_r(1)}$ for some computable function $f$. To overcome this barrier, we non-trivially use the underlying structure of $r$-list satisfying assignments. We further note that our reduction runs in time $|X|^{O_r(1)}$, which is polynomial also in $|X|$. Therefore, our methods give a \emph{unified proof of both the Baby PCP theorem and Baby PIH.} \footnote{Barto and Kozik \cite{BK22} derive Baby PCP with the stronger projection property. Using our techniques, we can get Baby PCP or Baby PIH with rectangular constraints, which is a slightly weaker property but still enough for many downstream reductions.}

As we mentioned earlier, PIH implies Baby PIH, and thus our result can be viewed as a first step towards proving the former. An intermediate problem between them is the following average version of Baby PIH: let an $r$-average list assignment be a labeling $L(u)$ for each variable $u$ such that the average cardinality of $L(u)$ over all the variables $u$ is at most $r$. 
\begin{conjecture}[Average Baby PIH]
For any constants $r >1,C$ and any computable function $f$, no algorithm can given as input a 2CSP instance on $k$ variables with size $n$, distinguish between the following two cases in $f(k)\cdot n^{C}$ time:
    \begin{itemize}
        \item (Completeness) There is an assignment satisfying all the constraints.
        \item (Soundness) No $r$-average list assignment satisfies all the constraints.
    \end{itemize}
\end{conjecture}

Note that the difference between the Baby and Average Baby versions is to use the $\ell_\infty$ vs. $\ell_1$ norms of the number of values given to the variables.
Once again, Average Baby PIH has a counterpart in the PCP world, namely the minimization version of Label Cover with $\ell_1$ total cost, which fueled early inapproximability results for {\sc SetCover} and basic problems that concern linear codes and lattices~\cite{ABSS97}, as surveyed in \cite{AL96}. 

Average Baby PIH is an intriguing open problem and could help in making further progress towards PIH. Furthermore, the Average Baby PIH is strong enough to imply some non-trivial inapproximability results. 
Notably, with an additional property called rectangular constraints (which we will define later), it implies constant inapproximability of \kexact{}, which we previously only knew under PIH \cite{Man20}. 

Towards a better understanding of Average Baby PIH, we give a counterexample showing that the direct product construction we use to prove Baby PIH as is cannot establish the average version. This suggests in order to get Average Baby PIH or full PIH, one may need other techniques or constructions. As a candidate, we pose a question that whether the $\wone$-hardness of approximating \kclique{} can help us bypass this counterexample. Please refer to \Cref{sec:discuss} for details. 

\medskip \noindent \textbf{Organization.}
In Section \ref{sec:pre}, we introduce some preliminaries, including the problems considered in this paper and related complexity hypotheses. In Section \ref{sec:baby}, we prove Baby PIH via the direct product construction. We then discuss Average Baby PIH in Section \ref{sec:avg}, and conclude with some open problems in \Cref{sec:discuss}.

\iffalse
\medskip \noindent \textbf{Paper Outline}. 
\begin{itemize}
    \item approximation, parameterization as two ways to cope, PIH
    \item list version of PIH
    \item Highlight baby PCP 
    \item our proof 
    \item Highlight why Barto-Kozik won't work immediately.
    \item we develop a simplified (than baby PCP) version of  baby PIH by requiring just rectangular constraints (which suffice for many reductions) compared to projection constraints of baby PCP. Comment that our reduction has poly k time 
    \item Progress towards baby PIH
    \item l1 vs linf
    \item baby PIH vs baby PCP. can we comment on why baby PIH doesn't imply baby PCP? and about having no implications of PIH in the PCP world?
    \item Add citations
\end{itemize}
\fi